\section{Products, convolution, and quantitative limit theorems}
\label{sec:limit-theorems}

\subsection{Independence and product integration}

Independence is supplied as a product construction, or by a proved product
factorization on a generating family; it is not inferred from a drawing or
from uncorrelatedness.  For finite rational laws, product weights are
$w_i v_j$ and both Fubini and
$\mathbb E[f(X)g(Y)]=\mathbb E f(X)\,\mathbb E g(Y)$ are identities of
finite sums.  For represented laws, use rational product approximants and
explicit absolute-integrability tails.  The product error estimate
\[
 |uv-u'v'|\leq |u|\,|v-v'|+|v'|\,|u-u'|
\]
and truncation reduce the statement to the finite identity.  A new
integrability condition is needed when the observable is unbounded.

If independent $X,Y$ have represented densities $f,g$, then their sum has
convolution density
\[
 h(z)=\int_{\mathbb R}f(z-y)g(y)\,dy,
\]
provided the density computation has local moduli and uniform tails on the
requested compact $z$-range.  This condition is automatic for the four
explicit families above and for a bounded uniformly continuous $f$ with a
known modulus and an effectively tight $g$: the discarded integral is at
most $\|f\|_\infty$ times the tail mass of $g$.  Normalization and the
pushforward formula are finite double-sum identities followed by the
specified tail limits.  Moreover $h\leq\|f\|_\infty$.

For independent centered normal variables with scales $\sigma,\tau>0$,
completing the square gives
\[
 \frac{(z-y)^2}{\sigma^2}+\frac{y^2}{\tau^2}
 =\frac{z^2}{\sigma^2+\tau^2}
  +\frac{\sigma^2+\tau^2}{\sigma^2\tau^2}
       \left(y-\frac{\tau^2z}{\sigma^2+\tau^2}\right)^2.
\]
The affine substitution in the last Gaussian integral proves
$X+Y\sim N(0,\sigma^2+\tau^2)$.  Every tail needed here is bounded by
\eqref{eq:gaussian-tail}; no characteristic-function uniqueness theorem is
required.  Induction gives
$n^{-1/2}\sum_{i=1}^nG_i\sim N(0,1)$.

\subsection{A Gaussian conditional density}

A useful continuous conditioning example follows from the same algebra.
Let $X,G$ be independent standard normals, let $\sigma\geq\sigma_0>0$, and
put $Y=aX+\sigma G$.  Write $A=a^2+\sigma^2$.  The joint density is
\[
 f(x,y)=\frac1{\sigma Z^2}
  \exp\!\left[-\frac12\left(x^2+\frac{(y-ax)^2}{\sigma^2}\right)\right],
\]
and the marginal of $Y$ is
$m(y)=(\sqrt A Z)^{-1}e^{-y^2/(2A)}$.  Completing the square gives
\[
 x^2+\frac{(y-ax)^2}{\sigma^2}
 =\frac{A}{\sigma^2}\left(x-\frac{ay}{A}\right)^2+\frac{y^2}{A}.
\]
Thus the density kernel $f(x,y)/m(y)$ is a normal density with mean
$ay/A$ and variance $\sigma^2/A$.  It is a defined conditional-density
version at every represented $y$.  It is not formed by dividing by the
zero probability of the singleton event $\{Y=y\}$.

The denominator separation is explicit on compact parameter ranges.
If $A\leq U^2$ with rational $U\geq1$ and $|y|\leq R$, then
\[
 m(y)\geq
 \frac1{4U\,3^{\lceil R^2/(2\sigma_0^2)\rceil}}.
\]
The finite product identity and the same Gaussian tails establish the
conditional integration identity against bounded effectively integrable
test functions.  The lower bound displayed above gives the local precision
required for division and makes the usual Bayesian Gaussian update a
primitive recursive operation on the supplied presentations.

\subsection{Markov, Chebyshev, and the weak law}

For a nonnegative finite observable $Y$ and rational $a>0$,
$a\mathbf1_{\{Y\geq a\}}\leq Y$.  Multiplying by weights and summing gives
Markov's inequality.  For a centered $X$ this yields
\[
 \mathbb P\{|X|\geq\varepsilon\}
 \leq\frac{\mathbb E X^2}{\varepsilon^2}.
\]
For pairwise independent centered $X_1,\ldots,X_n$ with
$\mathbb E X_i^2\leq V$, finite expansion and vanishing cross terms give
$\mathbb E(\sum_iX_i)^2\leq nV$.  Consequently
\begin{equation}
 \mathbb P\!\left\{\left|\frac1n\sum_{i=1}^nX_i\right|
                \geq\varepsilon\right\}
 \leq\frac{V}{n\varepsilon^2}.
 \label{eq:weak-law}
\end{equation}
A requested error probability $\delta>0$ is achieved by
$n\geq\max(1,\lceil V/(\varepsilon^2\delta)\rceil)$.
For noncentered identically distributed variables apply this to
$X_i-\mu$ with a supplied mean and variance certificate.

For general $L^2$ names, a sharp threshold need not itself be an effective
event.  The universal conclusion is then a certified upper bound, obtained
from finite approximants with rational threshold slack and the $L^2$
approximation error.  For the continuous density laws above, sharp threshold
events are available: a bounded density controls boundary strips.  In a sum
of independent density variables, convolution preserves any one available
global density bound; scaling by $1/n$ multiplies that bound by $n$.
Thus every finite threshold or finite maximum event in this section has
an explicit approximation modulus for these laws.  This is precisely what
is needed to pass the finite inequalities to the represented variables.

\subsection{A strong law with an explicit finite-window modulus}

Let $X_i$ be independent, centered, and satisfy
$\mathbb E X_i^2\leq V$.  The input is a single primitive recursive family
of coordinate and finite-product presentations, with the needed moment and
boundary certificates uniform in the coordinate index and finite dimension.
Separate, nonuniform external descriptions of each $X_i$ would not suffice.
Write $S_j=\sum_{i=1}^jX_i$.
For every $m\geq1$ and $a>0$, the finite Kolmogorov maximal inequality is
\begin{equation}
 \mathbb P\{\max_{1\leq j\leq m}|S_j|\geq a\}
 \leq mV/a^2.
 \label{eq:kolmogorov}
\end{equation}
Here full independence, or the appropriate conditional-zero-increment
property, is used; pairwise independence alone does not justify the proof.

\begin{proof}
Let $A_j$ be the event that the threshold is first reached at time $j$.
The events are disjoint, their union is the maximum event, and on $A_j$,
$S_j^2\geq a^2$.  Since $A_j$ and $S_j$ depend only on the first $j$
coordinates, independence and centering give
$\mathbb E[\mathbf1_{A_j}S_j(S_m-S_j)]=0$.  Expansion therefore yields
\[
 \mathbb E[\mathbf1_{A_j}S_m^2]
 =\mathbb E[\mathbf1_{A_j}S_j^2]
  +\mathbb E[\mathbf1_{A_j}(S_m-S_j)^2]
 \geq a^2\mathbb P(A_j).
\]
Sum over $j$ and use $\sum_j\mathbf1_{A_j}\leq1$ and
$\mathbb E S_m^2\leq mV$.  On rational finite spaces this is only an
identity of weighted sums plus positivity.  Moment and boundary certificates
justify the same assertion for the continuous laws.
\end{proof}

For any integers $1\leq N\leq M$, split $[N,M]$ into the dyadic blocks
$[2^jN,2^{j+1}N)$.  Failure of $|S_n|/n<\varepsilon$ in block $j$ implies
$\max_{n\leq2^{j+1}N}|S_n|\geq2^jN\varepsilon$.  By
\eqref{eq:kolmogorov} its probability is at most
$2^{1-j}V/(N\varepsilon^2)$.  The finite geometric sum is less than $4$;
therefore
\begin{equation}
 \mathbb P\!\left\{\max_{N\leq n\leq M}\frac{|S_n|}{n}
                         \geq\varepsilon\right\}
 \leq\frac{4V}{N\varepsilon^2}.
 \label{eq:strong-window}
\end{equation}
In particular, for $\varepsilon=2^{-r}$ and error probability $2^{-k}$,
\begin{equation}
 N(r,k)=\max\!\left(1,\left\lceil4V\,2^{2r+k}\right\rceil\right)
 \label{eq:strong-modulus}
\end{equation}
works \emph{uniformly in every finite endpoint $M\geq N(r,k)$}.
This is a substantive finite, primitive recursive form of the strong law.

When each finite deviation event has a uniform effective event name,
there is also a precise almost-sure certificate in the completed event
algebra.  Put $D_n^\varepsilon=\{|S_n|/n\geq\varepsilon\}$ and let
$F_{N,T}^\varepsilon$ be its union for $N\leq n\leq T$.  For $T'\geq T$,
\eqref{eq:strong-window} bounds the distance between the two finite unions by
\[
 d(F_{N,T'}^\varepsilon,F_{N,T}^\varepsilon)
 \leq\frac{4V}{(T+1)\varepsilon^2}.
\]
Thus their event-metric limit $A_N^\varepsilon$ exists with an explicit
primitive recursive modulus, and
$\mathbb P(A_N^\varepsilon)\leq4V/(N\varepsilon^2)$.  This supplies a
countable union by its residual-event bound; it does not incorrectly sum
the generally divergent individual Chebyshev estimates.

For each error parameter $p$, define the admitted union
\[
 U_p=\bigcup_{r\geq0}A_{N(r,p+r+1)}^{\,2^{-r}}.
\]
The $r$th event has mass at most $2^{-p-r-1}$, so this union has a geometric
residual bound and $\mathbb P(U_p)\leq2^{-p}$.  For every $r$ and every
$n\geq N(r,p+r+1)$, the deviation event $D_n^{2^{-r}}$ is included in
$U_p$, in the order of the event algebra.  This is an explicit almost-sure
convergence certificate: one small exceptional event simultaneously
controls all accuracy scales, with a supplied modulus outside it.
Under an ordinary set realization it implies the usual strong law.
The first-order content is the uniform finite-window estimate, the
explicitly coded exceptional events, and these inclusion estimates; it is
not quantification over an unrepresented powerset of all infinite sample
sequences.  For general $L^2$ names without threshold-event certificates,
retain the upper-envelope finite-window conclusion rather than claiming
these additional event names.

\subsection{An elementary quantitative central limit theorem}

The following rate is intentionally weaker than the optimal classical
Berry--Esseen rate; its advantage is a short proof with completely explicit
primitive recursive bounds.

Let $X_1,\ldots,X_n$ be independent and identically distributed, with
\[
 \mathbb E X_i=0,\qquad \mathbb E X_i^2=1,\qquad
 \mathbb E|X_i|^3\leq B,
\]
where $B$ is a supplied rational bound and the third absolute moment has
an effective tail certificate.  Let $G_i$ be independent standard normal
variables, independent also of all the $X_i$.  If a represented test
function $h$ is twice continuously differentiable, with second derivative
$L$-Lipschitz, then
\begin{equation}
 \left|\mathbb E h\!\left(n^{-1/2}\sum_{i=1}^nX_i\right)
               -\mathbb E h(G)\right|
 \leq\frac{L(B+2)}{6\sqrt n},
 \label{eq:smooth-clt}
\end{equation}
provided the displayed expectations are effectively integrable.  Bounded
$h,h',h''$, as in the application below, suffice.

\begin{proof}
The integral Taylor remainder gives, for all $x,u$,
\[
 h(x+u)=h(x)+h'(x)u+\tfrac12h''(x)u^2+R(x,u),\qquad
 |R(x,u)|\leq L|u|^3/6.
\]
This follows directly by integrating
$|h''(x+t)-h''(x)|\leq L|t|$ twice; the calculation is valid for negative
$u$ after reversing the integration limits.  It can alternatively be
proved by rational Riemann sums with the same error bound.
Replace the summands one at a time.  For replacement $i$, let $T_i$ be the
sum of the other normalized coordinates, using already replaced Gaussian
coordinates and the still unreplaced $X$ coordinates.  It is independent
of both $X_i$ and $G_i$.  Conditional on a fixed value of $T_i$, the constant,
linear, and quadratic Taylor terms have equal expectations for
$u=X_i/\sqrt n$ and $u=G_i/\sqrt n$, since their first two moments match.
The two remainder expectations have total magnitude at most
\[
 \frac{L}{6n^{3/2}}
      \bigl(\mathbb E|X_i|^3+\mathbb E|G_i|^3\bigr)
 \leq\frac{L(B+2)}{6n^{3/2}}.
\]
Product integration justifies averaging this bound in $T_i$.  Sum the $n$
replacement errors and use the proved Gaussian convolution identity.
For a direct arithmetic implementation, first truncate and approximate all
coordinates with the supplied third-moment tails, prove the finite
replacement bound with its recorded moment errors, and allocate the
requested extra precision across the $n$ replacements.  The exact matching
identities are recovered by those explicit error bounds, not by silently
asserting that rounded finite approximants have exactly unit variance.
\end{proof}

To turn this into a distribution-function theorem, use the explicit cutoff
\[
 H(u)=\begin{cases}
 1,&u\leq0,\\
 1-10u^3+15u^4-6u^5,&0\leq u\leq1,\\
 0,&u\geq1.
 \end{cases}
\]
It satisfies $0\leq H\leq1$, because
$H'(u)=-30u^2(1-u)^2\leq0$ in the transition interval.
The function is $C^2$, and $H''$ is globally $60$-Lipschitz: its third
derivative in $(0,1)$ is $-60+360u-360u^2$, which lies between $-60$ and
$30$, while its second derivative vanishes at both endpoints and outside
the interval.  No assertion of a globally continuous third derivative is
needed.

For $h_{t,\eta}(x)=H((x-t)/\eta)$ the Lipschitz constant is $60\eta^{-3}$,
and
\[
 \mathbf1_{(-\infty,t]}\leq h_{t,\eta}
 \leq\mathbf1_{(-\infty,t+\eta]}.
\]
The analogous cutoff starting at $t-\eta$ gives the lower bracket.  Since
$\varphi\leq1$, the Gaussian mass in an interval of length $\eta$ is at most
$\eta$.  Applying \eqref{eq:smooth-clt} to both brackets gives
\begin{equation}
 \sup_{t\in\mathbb R}|F_n(t)-\Phi(t)|
 \leq\eta+\frac{10(B+2)}{\sqrt n\,\eta^3}.
 \label{eq:cdf-clt}
\end{equation}
The supremum notation abbreviates the uniform inequality for every
represented real $t$; no operation computing an arbitrary supremum is
required.  For a requested error $\delta=2^{-k}$, set $\eta=\delta/2$.  It is
enough to take
\begin{equation}
 n\geq N_{\rm CLT}(k)
 :=\left\lceil160(B+2)\right\rceil^2\,2^{8k}.
 \label{eq:clt-rate}
\end{equation}
More generally the smooth theorem holds without identical distributions if
each summand has mean zero, variance one, and the same third-moment bound.

For continuous summand laws with the boundary certificates already given,
$F_n$ is an effective distribution function, so \eqref{eq:cdf-clt} is an
ordinary quantitative CDF assertion.  Without that event information, the
same proof supplies the upper and lower smooth-probability brackets.  It
does not manufacture a CDF name at a threshold whose boundary mass lacks the
required certificate.  For an original
law with mean $\mu$ and variance $\sigma^2>0$, supply a rational lower bound
on $\sigma$ and apply the theorem to $(X_i-\mu)/\sigma$.

\subsection{What remains to be formalized for these results}

The finite arithmetic is not starting from zero.  In the inspected ProveIt
snapshot, the directory
\texttt{Derivations/Lean/IPRADerivations/Probability/} contains object-level
\texttt{PRAAx} derivations for finite probability.  In particular,
\texttt{WeakLaw.lean} contains \texttt{weak\_law},
\texttt{FirstCrossing.lean} contains \texttt{cover}, \texttt{disjoint}, and
\texttt{crossing\_bound}, and the \texttt{Kolmogorov*} files assemble the
finite maximal argument.  These are relevant providers, not merely
mathematical prose.

The limits of that reuse matter.  \texttt{NormalMoments.lean} proves the
arithmetic of double factorials and explicitly says that the integrals
identifying them with normal moments are not formalized.
\texttt{CompleteSquare.lean} proves finite polynomial identities and
explicitly excludes normalization integrals.  \texttt{ProductReplacement.lean}
proves a finite product telescope and bounds, while its header excludes
characteristic functions and smoothing.  The Lindeberg proof here uses a
sum of expectation differences rather than that characteristic-function
product theorem, though the finite telescope infrastructure remains useful.

A completion proceeds in the following concrete order.  First define
rational quadrature and its finite error predicate, then compile the
exponential series estimates and the kernel-tail contracts.  Next construct
normalization and reciprocal names with their separation constants, and
prove moment identities by finite integration-by-parts estimates.  Prove
product integration with explicit truncation errors, boundary-strip transfer,
and Gaussian convolution.  The finite weak and maximal inequalities can
then be transported to the represented variables.  Finally encode the
quintic cutoff, its $60$-Lipschitz second-derivative certificate, the Taylor
remainder estimate, and the one-coordinate replacement inequality.  A
quantifier-free induction on the number of replacements proves the finite
sum bound.  Every precision choice in this section is an explicit primitive
recursive expression or a bounded search with a displayed bound.
